% =====================================================================
% MoVoC: Morphology-Aware Subword Construction for Ge'ez Script Languages
% CORRECTED VERSION of Findings of EMNLP 2025, pp. 13131--13144
% https://aclanthology.org/2025.findings-emnlp.706/
%
% Compile with XeLaTeX (Ethiopic script):  latexmk -xelatex main.tex
%
% Every result in this paper comes from github.com/hailaykidu/MoVoC.
% Tables are \input from tables/, produced by
%   python scripts/make_correction_tables.py --repo <MoVoC> --out tables
% (tables 2-4 copied verbatim from the repository; 1, 5, 6, 7 generated
% from repository data). No result value is typed by hand in this file;
% numbers quoted in the prose are checked by scripts/check_numbers.py.
% =====================================================================
\documentclass[11pt]{article}
\usepackage[final]{acl}

\usepackage{fontspec}
\setmainfont{TeX Gyre Termes}
\newfontfamily\ethiopicfont{NotoSerifEthiopic}[
  Path=fonts/, Extension=.ttf,
  UprightFont=*-Regular, BoldFont=*-Bold, ItalicFont=*-Regular, Scale=0.92]
\newcommand{\eth}[1]{{\ethiopicfont #1}}

\usepackage{microtype}
\usepackage{booktabs}
\usepackage{amsmath}
\usepackage{algorithm}
\usepackage{algpseudocode}
\usepackage{graphicx}
\usepackage{xcolor}

\newcommand{\movoctok}{MoVoC-Tok}

\title{MoVoC: Morphology-Aware Subword Construction for\\ Ge'ez Script Languages\\[2pt]
\normalsize\normalfont Corrected version}

\author{Hailay Kidu Teklehaymanot \and Dren Fazlija \and Wolfgang Nejdl \\
  L3S Research Center, Leibniz University Hannover, Germany \\
  \texttt{teklehaymanot@l3s.de}}

\begin{document}
\maketitle

\begin{abstract}
Subword-based tokenization methods often fail to preserve morphological
boundaries, a limitation especially pronounced in low-resource, morphologically
complex languages such as those written in the Ge'ez script. To address this,
we present Morpheme-aware Subword Vocabulary Construction (MoVoC) and train
\movoctok{}, a tokenizer that integrates supervised morphological analysis into
the subword vocabulary. This hybrid segmentation approach combines
morpheme-based and Byte Pair Encoding (BPE) tokens to preserve morphological
integrity while maintaining lexical meaning. To tackle resource scarcity, we
release morpheme-annotated data for four Ge'ez script languages (169,806
records) and a morpheme-aware vocabulary for two of them. Scoring every
comparison on words held out from the morpheme resource that constrains the
\movoctok{} applied to that language, \movoctok{} has the highest MorphScore
in three of four languages and the highest morpheme boundary precision in two
of four, with its largest and most clearly resolved gains on Tigre, a language
it was not trained on and to which it is applied cross-lingually. The machine
translation results reported in the published version are withdrawn. Data,
tokenizers and evaluation code are publicly
available.\footnote{\url{https://github.com/hailaykidu/MoVoC}}
\end{abstract}

% ---------------------------------------------------------------------
\section*{Correction Notice}
\label{sec:notice}

This is a corrected version of \citet{teklehaymanot-etal-2025-movoc}. The
evaluation artifacts (checkpoints, predictions, scoring pipeline) behind the
results in the published version were not recovered, and the published figures
could not be reproduced from the released code and data. Following the advice
of the EMNLP~2025 programme chairs, every result in this version is replaced
by results reconstructed from the public repository. The method
(Section~\ref{sec:method}) is unchanged. The main corrections are:

\begin{itemize}\setlength\itemsep{0pt}
\item \textbf{Evaluation sets (Section~\ref{sec:heldout}).} Part of each
  evaluation set consisted of words whose morpheme boundaries had constrained
  the training of the \movoctok{} that scores it, which the baselines never
  saw. All results are now scored on held-out words only: 76,488 of 81,224
  Amharic, 2,438 of 5,224 Tigrinya, 1,787 of 1,974 Tigre and 170 of 172 Ge'ez
  words. Tigre and Ge'ez are affected because the Tigrinya \movoctok{} scores
  them cross-lingually and the languages share surface forms.
\item \textbf{Intrinsic results (Tables~\ref{tab:morphscore} and~\ref{tab:intrinsic}).}
  On held-out words \movoctok{} leads on MorphScore in three of four languages
  and on boundary precision in two of four, not all four. Its largest gains are
  on Tigre, which is evaluated cross-lingually. All values are lower than
  published (e.g.\ precision 0.24--0.56 rather than 73.9--88.3). The published
  Table~2 ``No.~Items'' column gave pooled resource sizes, not evaluation-set
  sizes.
\item \textbf{Translation results.} The published Table~3 is withdrawn. The
  decode outputs are not in the repository, the recorded hypothesis/reference
  length ratios (2.06--15.10) show generation failing to terminate, and the
  decoding configuration and model initialisation of those runs cannot be
  established from the surviving records. This version makes no extrinsic
  claim.
\item \textbf{Examples (Tables~\ref{tab:segmentation} and~\ref{tab:annotation}).}
  The released tokenizers do not produce the segmentations shown in the
  published Table~1 and Section~7, and the published translation example could
  not be reproduced. Both are replaced by actual tokenizer output.
\item \textbf{Setup.} Training hyperparameters, evaluation data and vocabulary
  sizes are corrected to those recorded in the repository.
\end{itemize}
Appendix~\ref{app:corrections} lists every change.

% ---------------------------------------------------------------------
\section{Introduction}
\label{sec:intro}

Tokenization is a fundamental preprocessing step that converts raw text into
structured units such as bytes, characters, subwords, words or multi-word
expressions. Subword methods such as BPE \citep{sennrich-etal-2016-neural}
became popular because they are language-independent and compress the
vocabulary, allowing efficient learning of token representations.

Subwords, however, often fail to capture morphological structure. This is
especially problematic in multilingual models with a shared vocabulary, where
low-resource languages typically receive fewer subwords and therefore high
token-to-word ratios. Morphological complexity is associated with worse
language-model performance \citep{park-etal-2021-morphology,
arnett-bergen-2025-language}: morphologically rich languages produce many
infrequent word forms from combinations of morphemes, causing data sparsity,
and they tend to have less annotated data and non-Latin scripts that require
more UTF-8 bytes. Compression alone is therefore not always predictive of
downstream success.

Morphologically aware tokenization yields more meaningful tokens, and such
efforts have been constrained by the limited availability of morphologically
annotated data, mostly for high-resource languages. Relying solely on
morpheme-based tokenization is also suboptimal; hybrid methods that integrate
morphological information into subword tokenization share statistical
information across related word forms while keeping sequences short.

\paragraph{Contributions.}
(i)~We release morphologically annotated datasets for four low-resource Ge'ez
script languages to support tokenization, morphological analysis and
downstream NLP tasks. (ii)~We propose MoVoC, a supervised approach that
integrates a BPE-derived vocabulary with morphemes obtained from supervised
morphological analysis, together with \movoctok{}, a BPE tokenizer whose
merges may not cross morpheme boundaries. (iii)~We evaluate the approach
intrinsically, on words held out from the morpheme resource that constrained
training, and report where it helps and where it does not.

% ---------------------------------------------------------------------
\section{Background}
\label{sec:background}

\subsection{Tokenization Approaches}

\paragraph{Word tokenization} divides text into words, in its simplest form by
whitespace, or by rule-based patterns such as regular expressions. Its main
weakness is out-of-vocabulary (OOV) words arising from typos and unseen forms.

\paragraph{Subword tokenization} decomposes words into subwords, so that OOV
words can be composed from vocabulary units. Widely used algorithms include
WordPiece \citep{schuster2012japanese}, BPE \citep{sennrich-etal-2016-neural},
byte-level BPE and the Unigram model of SentencePiece
\citep{kudo-richardson-2018-sentencepiece}.

\paragraph{Character tokenization} at the character or UTF-8 byte level
increases sequence length, and with it the cost of self-attention, which is
quadratic in sequence length.

\subsection{Ge'ez Script Natural Language Processing}

Ge'ez is one of the oldest continuously used scripts, preserved for over two
thousand years. Unlike alphabetic systems such as Greek or Latin, it is a
syllabic script (an abugida) originally derived from the Sabaean script. It is
used for several morphologically rich Semitic languages, including Ge'ez,
Tigrinya, Amharic, Tigre and Blin. Verbs inflect for tense, aspect, person,
number and gender: the Amharic verb \eth{ጻፈ} (`he wrote') yields \eth{ጻፍኩ} (`I wrote'),
\eth{ጻፍክ} (`you wrote') and \eth{ጻፉ} (`they wrote'), and inflectional endings attach
directly to the stem without clear word boundaries. The development of usable
NLP tools for these languages is hindered by the scarcity of essential
linguistic resources. We contribute annotated data for four Ge'ez script
languages (Tigrinya, Amharic, Ge'ez and Tigre) and a morpheme-aware
alternative to BPE that constrains segmentation to respect morpheme
boundaries.

\subsection{Tokenization for Morphologically Rich and Low-Resource Languages}

\citet{park-etal-2021-morphology} trained language models with several
segmentation algorithms on Bible verses in 92 languages and found that models
perform worse on morphologically complex languages, and that linguistically
motivated segmentation reduces this effect. MorphBPE
\citep{asgari2025morphbpe} and MorphPiece \citep{jabbar2023morphpiece}
incorporate linguistic features into tokenization for morphologically rich
languages, and \citet{bauwens-delobelle-2024-bpe} argue that a lack of
morpheme awareness leads to inconsistent intra-word representations and
inflated vocabularies. Results are mixed, however:
\citet{saleva-lignos-2021-effectiveness} found that morphology-based
segmentation for low-resource NMT occasionally outperforms BPE but is often
statistically similar, and \citet{libovicky-helcl-2024-lexically} report that
semantically informed segmentation helps less for low-resource languages.
Our corrected results (Section~\ref{sec:results}) are consistent with this
mixed picture.

% ---------------------------------------------------------------------
\section{Proposed Method}
\label{sec:method}

The method has three parts: (1)~pre-tokenization and supervised morphological
analysis, (2)~vocabulary construction (MoVoC) and (3)~morpheme-aware subword
segmentation (\movoctok{}).

\subsection{Pre-tokenization and Supervised Morphological Analysis}
\label{sec:pretok}

Ge'ez-script languages lack robust NLP tools for basic preprocessing. We use a
pre-tokenization pipeline of regular expressions tailored to the orthography,
followed by corpus cleaning and supervised morphological analysis. For Amharic
and Tigrinya we use HornMorpho \citep{gasser2011hornmorpho}, a rule-based
morphological analyzer, with human post-editing; Tigrinya output in particular
often requires manual correction because of limited coverage. For Ge'ez and
Tigre we annotate morphemes manually.

Frequency-based BPE is unaware of morpheme boundaries and can split words
inside morphemes. Table~\ref{tab:segmentation} shows the segmentations that
the released BPE, WordPiece and \movoctok{} tokenizers actually produce for
the example words used in the published version of this paper. None of the
three tokenizers consistently recovers the morpheme structure of these words;
for \eth{መምህርነት} (`profession of teaching') \movoctok{} and BPE produce
the same split, and only WordPiece separates the nominalizer \eth{-ነት}.
These examples are illustrative only; Section~\ref{sec:results} gives the
aggregate measurements.

\begin{table*}[t]
  \centering
  % Generated by make_correction_tables.py -- do not edit by hand.
\begin{tabular}{llllll}
\toprule
Word & Lang. & Gold & BPE & WordPiece & MoVoC-Tok \\
\midrule
\eth{አልሰበሩም} & Amharic & -- & \eth{አልሰ-በሩ-ም} & \eth{አልሰ-በሩ-ም} & \eth{አል-ሰበ-ሩም} \\
\eth{ኣይትከውንን} & Tigrinya & -- & \eth{ኣይት-ከውን-ን} & \eth{ኣይት-ከውን-ን} & \eth{ኣይት-ከውንን} \\
\eth{መምህርነት} & Amharic & \eth{መ-ም-ህርነት} & \eth{መምህ-ርነት} & \eth{መምህር-ነት} & \eth{መምህ-ርነት} \\
\bottomrule
\end{tabular}

  \caption{Segmentations produced by the released tokenizers (32k
  vocabularies) for the example words of the published Table~1 and
  Section~7. Amharic words use the Amharic \movoctok{}, the Tigrinya word the
  Tigrinya \movoctok{}. ``Gold'' is the surface segmentation in the evaluation
  set where the word occurs in it (``--'': not in the evaluation set).
  Generated by \texttt{make\_correction\_tables.py}.}
  \label{tab:segmentation}
\end{table*}

The annotated morphemes also serve as gold-standard test sets for evaluating
token-to-morpheme boundary alignment.

\subsection{Vocabulary Construction (MoVoC)}
\label{sec:movoc}

Standard subword methods select units by frequency (BPE) or likelihood
(Unigram) without considering morphemes. We design a hybrid vocabulary
construction (Algorithm~\ref{alg:movoc}) that allocates a portion of the
vocabulary budget to morpheme units and the rest to BPE tokens.

Let $P_{am}$ and $P_{ti}$ be the Amharic and Tigrinya monolingual corpora,
$s$ the total vocabulary budget and $r \in [0,1]$ the proportion reserved for
morphemes. Each language receives $s_{lang} = s/2$, of which
$s_{BPE} = s_{lang}(1-r)$ goes to a BPE model trained on that language and
$s_{morpheme} = s_{lang}\,r$ to morphemes.
\texttt{extract\_morphemes}$(P, s_{morpheme})$ segments the corpus with
HornMorpho, counts morpheme frequencies and keeps the
$k = s_{morpheme}$ most frequent: $V_{morpheme} = \mathrm{Top}_k(\mathrm{freq})$.
The final vocabulary is the union of the four sets. In the released
configuration $s = 224{,}000$ and $r = 5/7$, giving $s_{BPE} = 32{,}000$ and a
morpheme budget of $80{,}000$ per language. Fewer distinct morphemes were
found than the budget allows, so the realised vocabulary has 114,553 entries
(Table~\ref{tab:vocab}, Appendix~\ref{app:details}).

\begin{algorithm}[t]
\caption{MoVoC vocabulary construction}
\label{alg:movoc}
\small
\begin{algorithmic}[1]
\Require corpora $P_{am}$, $P_{ti}$; budget $s$; morpheme ratio $r$
\Ensure vocabulary $V_{MoVoC}$
\State $M_{am} \gets$ HornMorpho\_segment$(P_{am})$; $M_{ti} \gets$ HornMorpho\_segment$(P_{ti})$
\State $s_{lang} \gets s/2$; $s_{BPE} \gets s_{lang}(1-r)$; $s_{morpheme} \gets s_{lang}\, r$
\State $V_{BPE,am} \gets$ Train\_BPE$(P_{am}, s_{BPE})$; $V_{BPE,ti} \gets$ Train\_BPE$(P_{ti}, s_{BPE})$
\State $V_{morph,am} \gets$ extract\_morphemes$(P_{am}, s_{morpheme})$
\State $V_{morph,ti} \gets$ extract\_morphemes$(P_{ti}, s_{morpheme})$
\State $V_{MoVoC} \gets V_{BPE,am} \cup V_{BPE,ti} \cup V_{morph,am} \cup V_{morph,ti}$
\State \Return $V_{MoVoC}$
\end{algorithmic}
\end{algorithm}

\subsection{MoVoC-Tok: Morpheme-Aware Subword Segmentation}
\label{sec:movoctok}

A BPE tokenizer initialised with the MoVoC vocabulary can still learn merges
that cross morpheme boundaries. \movoctok{} therefore restricts merge
candidates during training. Let $W = \{w_1, \dots, w_n\}$ be the training words,
each a character sequence, and $M_i$ the morpheme boundaries of $w_i$. We seek
\begin{equation}
  \max_V \sum_{w_i \in W} \log P\bigl(\mathrm{BPE}(w_i; V, M_i)\bigr)
\end{equation}
subject to every unit $s_j$ of $\mathrm{BPE}(w_i; V, M_i) = (s_1, \dots, s_t)$
lying within $w_i$ without crossing a boundary in $M_i$: a merge $(a,b)$ is
permitted only if $ab$ does not span a boundary.

The released tokenizers apply this constraint to one language at a time with
a 32,000-entry vocabulary each, trained on 400,000 lines of NLLB text. Morpheme
boundaries come from 10,293 Amharic and 6,692 Tigrinya segmented records; the
constraint rejected 335,638 (Amharic) and 104,964 (Tigrinya)
boundary-crossing merges. At encoding time the tokenizer applies its learned
merges without access to gold boundaries.

% ---------------------------------------------------------------------
\section{Experimental Setup}
\label{sec:setup}

\subsection{Target Languages}

We study Amharic, Tigrinya, Ge'ez and Tigre. For Amharic and Tigrinya,
HornMorpho decomposes words into stems and affixes and the segmented units are
used for vocabulary construction. We also tried Morfessor
\citep{creutz2007unsupervised}; as an unsupervised model inferring boundaries
from surface patterns, it performed poorly. For Ge'ez and Tigre no analyzer or
training morpheme data was available; their manual annotations are used for
evaluation only, and the Tigrinya \movoctok{} is applied to them
cross-lingually.

\subsection{Datasets}
\label{sec:data}

\paragraph{Morpheme annotations.} The release contains 169,806 annotated
records: 153,759 Amharic (HornMorpho output with human post-editing), 7,531
Tigrinya post-edited records plus 206 Tigrinya gold records, 8,117 Tigre and
193 Ge'ez (both manually annotated). The intrinsic evaluation sets derived from
them contain 81,224 (Amharic), 5,224 (Tigrinya), 1,974 (Tigre) and 172
(Ge'ez) words, each with at least one gold morpheme boundary. All results are
reported on the held-out subsets of these sets (76,488, 2,438, 1,787 and 170
words; Section~\ref{sec:heldout} and Table~\ref{tab:morphscore}). The
published version reported 80k, 80k, 32k and 20k ``items''; those figures were
estimated pooled resource sizes, not the evaluation sets.

\subsection{Training Setup and Configuration}
\label{sec:config}

\paragraph{Tokenizers.} The BPE and WordPiece baselines are trained with the
Hugging Face \texttt{tokenizers} library (v0.21.4) to 32,000 entries on a
shared corpus of 2,796,346 sentences, with NFC normalisation and metaspace
pre-tokenization. These baselines and the per-language \movoctok{} tokenizers
(Section~\ref{sec:movoctok}) are used for all intrinsic evaluation.

% ---------------------------------------------------------------------
\section{Evaluation Framework}
\label{sec:eval}

\subsection{Intrinsic Evaluation}
\label{sec:intrinsic-def}

Gold boundaries are the character offsets between consecutive gold morphemes;
predicted boundaries are the offsets between consecutive tokens. All three
tokenizers are scored on identical word lists.

\paragraph{Morpheme boundary precision} is the fraction of all predicted
boundaries, pooled over the evaluation set, that coincide with a gold
boundary.

\paragraph{MorphScore} follows \citet{arnett-bergen-2025-language}: the
fraction of gold boundaries recovered, pooled over the words the tokenizer
segments. Words left unsegmented are excluded, and extra boundaries are not
penalised, so MorphScore is recall-oriented and complements precision.

\paragraph{Rényi entropy} \citep{renyi1961measures} at $\alpha = 2$,
$H_2 = -\log \sum_i p_i^2$, is computed over the distribution of tokens the
tokenizer produces on the evaluation words and reported in nats. Lower values
indicate a more concentrated token distribution.

\subsection{Held-out evaluation}
\label{sec:heldout}

\movoctok{} is a supervised tokenizer: it consumes morpheme annotations that
BPE and WordPiece do not receive. Using that supervision is the method, and a
gain obtained from it is a result. It is a result, however, only on words the
supervision was not given for. Part of each evaluation set is drawn from the
morpheme resource that constrained the training of the \movoctok{} that
scores it (Section~\ref{sec:movoctok}), and on those words the tokenizer
reproduces a
constraint it was handed rather than generalizing from it, while the baselines
meet the words for the first time.

We therefore exclude, per language, every evaluation word occurring in the
morpheme resource named as \texttt{morpheme\_source} in that tokenizer's
training configuration. This removes 4,736 of 81,224 Amharic words (5.8\%) and
2,786 of 5,224 Tigrinya words (53.3\%), leaving 76,488 and 2,438. Exclusion is
by set membership against the resource itself rather than by the evaluation
set's provenance labels, which record only the first source each word was drawn
from and so undercount the overlap: 18 of the 160 Tigrinya words labelled gold
and held out are in fact present in the training resource.

Tigre and Ge'ez have no training morpheme resource of their own, but the
tokenizer that scores them does. They are evaluated cross-lingually with the
Tigrinya \movoctok{}, which was constrained by
\texttt{tigrinya\_morpheme\_segmented.json}, and the three languages are
Ethio-Semitic and share surface forms, so the same exclusion applies: 187 of
1,974 Tigre words (9.5\%) and 2 of 172 Ge'ez words occur in that resource,
leaving 1,787 and 170. The resource is therefore selected by the tokenizer
that scores a language, not by the name of the language itself.

All results in Section~\ref{sec:results} are reported on held-out words.
Differences are assessed with a paired bootstrap over words (10,000 resamples,
resampling index shared across tokenizers), reported in
Table~\ref{tab:significance}.

% ---------------------------------------------------------------------
\section{Results}
\label{sec:results}

\begin{table}[t]
  \centering
  \small
  \resizebox{\columnwidth}{!}{% Table 2: intrinsic MorphScore, held-out words only.
% Generated by scripts/evaluate_intrinsic.py --exclude-training.
% Words covered by MoVoC-Tok's training morpheme resource are excluded.
\begin{tabular}{llrr}
\toprule
Language (ISO 639-3) & Tokenizer & No. Items & MorphScore $\uparrow$ \\
\midrule
Amharic (amh) & BPE & 76,488 & 0.4127 \\
Amharic (amh) & WordPiece & 76,488 & 0.3872 \\
Amharic (amh) & MoVoC-Tok & 76,488 & 0.4148 \\
Tigrinya (tir) & BPE & 2,438 & 0.3648 \\
Tigrinya (tir) & WordPiece & 2,438 & 0.3613 \\
Tigrinya (tir) & MoVoC-Tok & 2,438 & 0.3833 \\
Tigre (tig) & BPE & 1,787 & 0.4987 \\
Tigre (tig) & WordPiece & 1,787 & 0.4776 \\
Tigre (tig) & MoVoC-Tok* & 1,787 & 0.5302 \\
Ge'ez (gez) & BPE & 170 & 0.6640 \\
Ge'ez (gez) & WordPiece & 170 & 0.6364 \\
Ge'ez (gez) & MoVoC-Tok* & 170 & 0.6534 \\
\bottomrule
\end{tabular}
}
  \caption{MorphScore of BPE, WordPiece and \movoctok{} on the morpheme
  evaluation sets, on words held out from \movoctok{}'s training morpheme
  resource (Section~\ref{sec:heldout}; No.~Items = evaluation words). For
  Tigre and Ge'ez ($*$) the Tigrinya \movoctok{} is applied cross-lingually.}
  \label{tab:morphscore}
\end{table}

\begin{table}[t]
  \centering
  \small
  \resizebox{\columnwidth}{!}{% Significance of held-out differences, MoVoC-Tok minus baseline.
% Paired bootstrap over words, 10,000 resamples, shared resampling index.
% Generated by scripts/bootstrap_significance.py.
\begin{tabular}{llrrr}
\toprule
Language & Metric & $\Delta$ vs BPE & 95\% CI & $p$ \\
\midrule
Amharic & Precision & +0.0018 & [+0.0004, +0.0033] & 0.012 \\
Amharic & MorphScore & +0.0021 & [+0.0001, +0.0042] & 0.045 \\
Tigrinya & Precision & -0.0016 & [-0.0106, +0.0075] & 0.744 \\
Tigrinya & MorphScore & +0.0185 & [+0.0031, +0.0345] & 0.017 \\
Tigre* & Precision & +0.0250 & [+0.0060, +0.0446] & 0.009 \\
Tigre* & MorphScore & +0.0274 & [+0.0096, +0.0458] & 0.003 \\
Ge'ez* & Precision & -0.0025 & [-0.0260, +0.0195] & 0.839 \\
Ge'ez* & MorphScore & -0.0105 & [-0.0383, +0.0157] & 0.442 \\
\bottomrule
\end{tabular}
}
  \caption{Paired bootstrap over held-out words (10,000 resamples, shared
  resampling index), \movoctok{} minus BPE. Intervals are 95\% percentile
  intervals of the difference. $*$ marks cross-lingual application of the
  Tigrinya \movoctok{}.}
  \label{tab:significance}
\end{table}

\paragraph{MorphScore (Table~\ref{tab:morphscore}).} On held-out words
\movoctok{} has the highest MorphScore of the three tokenizers in Amharic
(0.4148), Tigrinya (0.3833) and Tigre (0.5302); on Ge'ez BPE is ahead (0.6640
against 0.6534) on 170 words, which does not resolve the difference
(Table~\ref{tab:significance}). The published version reported 0.654--0.731
for \movoctok{} and stated that all four languages exceed the mean MorphScore
for fusional languages reported by \citet{arnett-bergen-2025-language}; that
comparison does not hold for the reconstructed values and is withdrawn.

\paragraph{Boundary precision (Table~\ref{tab:intrinsic}).} On held-out words
\movoctok{} has the highest precision in Amharic ($+0.0018$ over BPE,
$p = 0.012$) and Tigre ($+0.0286$, $p = 0.004$). On Tigrinya the difference
against BPE is $-0.0016$ with a 95\% interval of $[-0.0106, +0.0075]$: the
held-out data do not separate the two tokenizers, and we claim no precision
gain for Tigrinya. On Ge'ez BPE is marginally ahead ($-0.0026$, $p = 0.829$),
a difference that a 170-word test set cannot resolve. The in-language margins
over BPE are small; those over WordPiece are an order of magnitude larger
($+0.0176$ on Amharic). Precision is higher for Tigre for all three
tokenizers; Tigre's annotations are fully concatenative, so its morpheme
boundaries coincide with character offsets. The published version reported
precision of 73.9--88.3 and a \movoctok{} lead in all four languages, with the
largest gains on Tigre and Ge'ez; neither holds.

The Tigrinya precision result is the one conclusion that changes under the
held-out protocol. Over the full evaluation set, including the words
\movoctok{} was trained on, it leads (0.3242 against 0.3142 for BPE); over
held-out words alone it does not. We read this as a statement about annotation
coverage rather than about the method: the Tigrinya morpheme vocabulary
realises 7,125 types against 50,978 for Amharic (Table~\ref{tab:vocab}), from
6,692 training records against 10,293, so the constraint reaches too few word
types for the learned merges to transfer beyond the annotated words. That the
recall-oriented metric still improves ($+0.0185$, $p = 0.017$) while precision
does not is consistent with this reading. The clearest evidence that morpheme
supervision generalizes is Tigre, which has no training morpheme resource of
its own and carries the largest and most clearly resolved gains in the table
($+0.0286$ precision, $p = 0.004$; $+0.0314$ MorphScore, $p < 0.001$, against
BPE on 1,787 held-out words).

On the 187 excluded words, BPE places 10 boundaries (8 correct) and
\movoctok{} 22 (11 correct); the set is too small to interpret. Excluding
them raises the Tigre margin rather than lowering it
(Table~\ref{tab:significance}).

\paragraph{Rényi entropy (Table~\ref{tab:intrinsic}).} \movoctok{} has lower
entropy than BPE in Amharic, Tigrinya and Tigre and higher on Ge'ez, but
WordPiece has the lowest entropy in Amharic, Tigrinya and Tigre and BPE on
Ge'ez, so \movoctok{} does not lead on this metric in any language. The
published version reported normalized entropy (0.39--0.49) with \movoctok{}
best everywhere; the values here are raw and not comparable to it.

\begin{table}[t]
  \centering
  \small
  \resizebox{\columnwidth}{!}{% Table 4: boundary precision and Renyi entropy, held-out words only.
% Generated by scripts/evaluate_intrinsic.py --exclude-training.
\begin{tabular}{llrr}
\hline
Language & Tokenization & Precision $\uparrow$ & R\'enyi Entropy $\downarrow$ \\
\hline
Amharic & MoVoC-Tok & 0.3174 & 6.0292 \\
Amharic & BPE & 0.3156 & 6.2159 \\
Amharic & WordPiece & 0.2998 & 5.9725 \\
Tigrinya & MoVoC-Tok & 0.2385 & 5.9457 \\
Tigrinya & BPE & 0.2401 & 6.0529 \\
Tigrinya & WordPiece & 0.2401 & 5.9193 \\
Tigre & MoVoC-Tok* & 0.5642 & 5.2165 \\
Tigre & BPE & 0.5356 & 5.3060 \\
Tigre & WordPiece & 0.5116 & 4.9293 \\
Ge'ez & MoVoC-Tok* & 0.4282 & 3.9753 \\
Ge'ez & BPE & 0.4308 & 3.8629 \\
Ge'ez & WordPiece & 0.4182 & 3.9158 \\
\hline
\end{tabular}
}
  \caption{Morpheme boundary precision and Rényi entropy ($\alpha=2$, nats)
  for 32k vocabularies, on words held out from \movoctok{}'s training morpheme
  resource (Section~\ref{sec:heldout}). For Tigre and Ge'ez ($*$) the Tigrinya
  \movoctok{} is applied cross-lingually.}
  \label{tab:intrinsic}
\end{table}

\paragraph{Machine translation: withdrawn.} The published version reported
BLEU and chrF++ for English$\rightarrow$X translation with each tokenizer.
Those results are withdrawn, for three reasons. First, none of the 54 decode
passes
recorded for those runs has its hypothesis and reference files in the released
repository or its history, so the scores cannot be checked against the
translations they summarise. Second, the recorded hypothesis-to-reference
length ratios run 2.06--4.28 on FLORES-200 and as high as 15.10 on the smaller
OPUS sets: decoding is reaching the 128-token cap instead of emitting
end-of-sequence, which depresses BLEU independently of translation quality and
is not a property of the tokenizers under comparison. A systematic failure to
terminate makes the resulting scores uninformative about what they were meant
to measure. Third, the
surviving records disagree on decoding, one giving five-beam search and
another greedy decoding with a single beam, and whether the models were
initialised from a pretrained MarianMT checkpoint or trained from scratch is
unresolved; a corrected paper cannot publish a results table whose setup it
cannot describe.

We therefore make no extrinsic claim. MoVoC's contribution is a tokenization
method and an annotated dataset, and we evaluate it intrinsically. Whether
morphology-aware tokenization improves downstream translation for these
languages remains open, and answering it requires a corrected generation
pipeline rather than a re-scoring of these runs.

% ---------------------------------------------------------------------
\section{Qualitative Analysis}
\label{sec:qualitative}

The published version illustrated \movoctok{} with segmentations of
\eth{ኣይትከውንን} (`do not do it') and \eth{መምህርነት}, and with an
Amharic$\rightarrow$English translation of \eth{ቤቱን} \eth{አላየሁም} (`I did not
see the house'). The released tokenizers do not produce those segmentations
(Table~\ref{tab:segmentation}): \movoctok{} splits \eth{ኣይትከውንን} as
\eth{ኣይት-ከውንን} and \eth{መምህርነት} as \eth{መምህ-ርነት}, the same as
BPE. Both words of the translation example are single tokens for all three
tokenizers, and this version reports no translation results, so that example
is withdrawn.

At the level of individual words, then, the boundary constraint does not
reliably yield morpheme-aligned units. Its effect is visible only in the
aggregate, as the small precision and MorphScore margins in
Tables~\ref{tab:morphscore} and~\ref{tab:intrinsic}. This is consistent with
how the constraint is applied: it acts during training on the 3,372 Amharic
and 1,470 Tigrinya corpus words whose morpheme boundaries could be located,
while at encoding time the learned merges are applied without morphological
information.

% ---------------------------------------------------------------------
\section{Conclusion and Future Work}
\label{sec:conclusion}

We release morphologically annotated data for four Ge'ez script languages and
propose MoVoC, a morpheme-aware vocabulary construction, with \movoctok{}, a
BPE tokenizer whose merges respect morpheme boundaries. Scoring every
comparison on words held out from the morpheme resource that constrains the
tokenizer applied to that language, \movoctok{} attains the highest MorphScore in
three of four languages and the highest boundary precision in two of four, by
small margins in-language. Its largest gains are on Tigre, which it was not
trained on and to which it is applied cross-lingually, indicating that the
morpheme constraint produces segmentation aligning with morpheme structure
beyond the words the supervision covered. Where annotation coverage is thinner
the effect is weaker: on Tigrinya, whose morpheme vocabulary realises 7,125
types against 50,978 for Amharic, the MorphScore advantage persists on
held-out words while the precision advantage does not. The translation results
are withdrawn, so this version does not establish whether the intrinsic
advantage carries over to translation quality. Future work includes
in-language tokenizers for Tigre and Ge'ez, larger evaluation sets for them,
and translation experiments with a corrected generation pipeline.

% ---------------------------------------------------------------------
\section*{Limitations}

\paragraph{No downstream evaluation.} The translation results are withdrawn
(Section~\ref{sec:results}), so the method is assessed only by how well its
segmentations align with gold morpheme boundaries. Boundary precision and
MorphScore measure that alignment directly, but they do not establish that
better-aligned segmentation improves any downstream task, and we do not claim
it does.

\paragraph{Held-out evaluation and its cost.} All results are reported on
words held out from the morpheme resource that constrained the \movoctok{}
applied to that language (Section~\ref{sec:heldout}). We did not re-derive those
resources to enlarge the held-out sets, so Tigrinya retains 2,438 evaluation
words and its confidence intervals are correspondingly wide. On Tigrinya the
boundary precision advantage does not survive the exclusion, and we claim no
precision gain there. The Amharic margins over BPE, while statistically
resolved on 76,488 words, are small in absolute terms ($+0.0018$ precision,
$+0.0021$ MorphScore), and we do not claim they are of practical consequence
on their own.

\paragraph{Evaluation-set size.} The held-out evaluation sets differ by more
than two orders of magnitude (76,488 Amharic words against 170 Ge'ez words);
results for Ge'ez and Tigre carry much wider uncertainty.

\paragraph{Cross-lingual application.} No in-language \movoctok{} exists for
Tigre or Ge'ez; their results measure transfer from the Tigrinya tokenizer.

\paragraph{Published results.} The artifacts behind the published version
were not recovered, so its figures cannot be traced; the results here are new
measurements, not a reproduction.

\paragraph{Scope.} The annotated data and vocabularies cover a small set of
Ge'ez script languages, and the evidence here is intrinsic only. Whether the
approach helps on downstream tasks is untested in this version; see
Section~\ref{sec:results} on the withdrawal of the translation results.

\section*{Ethical Considerations and Resource Use}

We use publicly available data (NLLB, OPUS, FLORES-200, HornMT and the Mermru
English--Ge'ez corpus) and the HornMorpho analyzer, in line with their
licenses and intended research use. The morpheme annotations and vocabularies
we release are intended for research and are published under open licenses.
The original manuscript's readability was improved with ChatGPT paraphrasing
and language editing. This corrected version was prepared with the assistance
of an AI model (Claude) for drafting and consistency checking; all reported
numbers were verified against the repository.

\section*{Acknowledgments}

This research was supported by the German Academic Exchange Service (DAAD)
through the Hilde Domin Programme (funding no.\ 57615863).

\bibliography{custom}

\appendix

% ---------------------------------------------------------------------
\section{Additional Details}
\label{app:details}

\begin{table}[h]
  \centering
  \small
  % Generated by make_correction_tables.py from configs/movoc_config.json.
\begin{tabular}{lrr}
\toprule
Vocabulary & Budget & Realised size \\
\midrule
Amharic morphemes & 80,000 & 50,978 \\
Amharic BPE & 32,000 & 32,000 \\
Tigrinya morphemes & 80,000 & 7,125 \\
Tigrinya BPE & 32,000 & 32,000 \\
Bilingual MoVoC (union) & 224,000 & 114,553 \\
\bottomrule
\end{tabular}

  \caption{MoVoC vocabulary budget and realised sizes (released
  configuration: $s=224{,}000$, $r=5/7$). The union is smaller than the sum
  of its parts because the sets overlap. The published version reported a
  152k bilingual vocabulary.}
  \label{tab:vocab}
\end{table}

\paragraph{Sources of difficulty.} The main difficulty in processing Ge'ez
script languages lies less in the script itself than in their morphology.
Amharic and Tigrinya are fusional: a single word encodes tense, aspect, mood,
person, number, gender and voice through prefixes, suffixes and stem-internal
changes, and individual morphemes often carry several grammatical meanings, so
morpheme boundaries are not always clear-cut. For example, Amharic \eth{ጻፈ} (`he
wrote'), \eth{ጻፈች} (`she wrote') and \eth{ጻፉ} (`they wrote') differ in their endings,
while \eth{ይጻፋል} (`he is writing') adds a prefix and a suffix. Frequency-based
tokenizers are agnostic to this structure.

\begin{table*}[t]
  \centering
  \small
  % Generated by make_correction_tables.py -- records copied from data/annotations/.
\begin{tabular}{llccccc}
\toprule
Language & Word & Prefix & Root & Infix & Suffix & Clitic \\
\midrule
Amharic & \eth{መምህርነት} & \eth{መ} & \eth{ም-ህ-ር} & -- & \eth{ነት} & -- \\
Tigrinya & \eth{ዝወደቐ} & \eth{ዝ-} & \eth{ወደቐ} & -- & -- & -- \\
Ge'ez & \eth{ለንኡሳት} & \eth{ለ-} & \eth{ንኡስ} & -- & \eth{-ኣት} & -- \\
Tigre & \eth{ናዊ} & -- & \eth{ና} & -- & \eth{-ዊ} & -- \\
\bottomrule
\end{tabular}

  \caption{Annotation records from the released data. Amharic shows the
  published example; the published Tigrinya, Ge'ez and Tigre examples are not
  in the released files for those languages, so the first multi-morpheme
  record of each file is shown. ``--'': empty slot.}
  \label{tab:annotation}
\end{table*}

\section{Annotation Guidelines}

Annotators split each word into morphemes and label each with one of PREFIX,
ROOT, SUFFIX, INFIX or CLITIC (Table~\ref{tab:annotation}), consistently across
words and languages; slots that do not apply are left empty. The guidelines
cover Amharic, Tigrinya, Ge'ez and Tigre.

\section{List of Corrections}
\label{app:corrections}

Relative to \citet{teklehaymanot-etal-2025-movoc}. Section and table numbers
are those of the published version; where a part of this version has no
published counterpart, or the numbering differs, it is given as a reference to
the section here.
\begin{enumerate}\setlength\itemsep{0pt}
\item Abstract: ``consistent improvements in intrinsic metrics'' corrected to
  three of four languages on MorphScore and two of four on boundary precision,
  measured on held-out words; the translation statement is removed.
\item Table~1: replaced by actual tokenizer output (Table~\ref{tab:segmentation}).
\item Section~3.2: BPE share corrected to $s_{lang}(1-r)$; realised
  vocabulary sizes added.
\item Section~4.2: evaluation-set sizes corrected; the description of the
  translation data is removed with the translation results.
\item Section~4.3: the tokenizer configuration is corrected; the translation
  model configuration is removed with the translation results.
\item Section~5.2: metric definitions aligned with the released evaluation
  code; entropy now reported raw. These are Section~\ref{sec:intrinsic-def}
  here.
\item New in this version (Section~\ref{sec:heldout}): every result is scored
  on words held out from the morpheme resource that constrained the
  \movoctok{} applied to that language, which removes 4,736 Amharic, 2,786
  Tigrinya, 187 Tigre and 2 Ge'ez evaluation words. For Tigre and Ge'ez the
  resource is the Tigrinya one, since the Tigrinya \movoctok{} scores them
  cross-lingually.
\item Table~2: held-out evaluation-set sizes and reconstructed MorphScores for
  all three tokenizers; the comparison with the fusional-language mean is
  withdrawn.
\item Table~3: withdrawn, together with the OPUS appendix; this version
  reports no translation results.
\item Table~4: held-out precision and entropy for all three tokenizers;
  ``best in all four languages'' corrected to two of four on precision. A new
  table reports bootstrap confidence intervals for the differences against BPE.
\item Section~7: segmentation examples replaced by actual output; translation
  example withdrawn.
\item Sections~8 and Limitations: conclusions aligned with the corrected
  results.
\item Table~5: realised vocabulary sizes (114,553, not 152k).
\item Table~6: examples replaced by records present in the released data.
\end{enumerate}

\end{document}
